\section{由流动场构造的能动量张量}
\label{sec:4}
\subsection{小流时间展开与重正化群论证}
为了把式~\eqref{eq:(2.7)} 和~\eqref{eq:(2.10)} 中的能动量张量用流动场表出，我们考虑由
流动场构成的下列二阶对称张量（它们在 CP 变换下为偶）：
\begin{align}
   \Tilde{\mathcal{O}}_{1\mu\nu}(t,x)&\equiv
   G_{\mu\rho}^a(t,x)G_{\nu\rho}^a(t,x),
\label{eq:(4.1)}\\
   \Tilde{\mathcal{O}}_{2\mu\nu}(t,x)&\equiv
   \delta_{\mu\nu}G_{\rho\sigma}^a(t,x)G_{\rho\sigma}^a(t,x),
\label{eq:(4.2)}\\
   \Tilde{\mathcal{O}}_{3\mu\nu}(t,x)&\equiv
   \mathring{\Bar{\chi}}(t,x)
   \left(\gamma_\mu\overleftrightarrow{D}_\nu
   +\gamma_\nu\overleftrightarrow{D}_\mu\right)
   \mathring{\chi}(t,x),
\label{eq:(4.3)}\\
   \Tilde{\mathcal{O}}_{4\mu\nu}(t,x)&\equiv
   \delta_{\mu\nu}
   \mathring{\Bar{\chi}}(t,x)
   \overleftrightarrow{\Slash{D}}
   \mathring{\chi}(t,x),
\label{eq:(4.4)}\\
   \Tilde{\mathcal{O}}_{5\mu\nu}(t,x)&\equiv
   \delta_{\mu\nu}
   m\mathring{\Bar{\chi}}(t,x)
   \mathring{\chi}(t,x).
\label{eq:(4.5)}
\end{align}
注意上述所有算符 $\Tilde{\mathcal{O}}_{i\mu\nu}(t,x)$ 对任意~$D$ 的量纲均为~$4$。

我们还在 $D$ 维 $x$ 空间引入相应的裸算符：
\begin{align}
   \mathcal{O}_{1\mu\nu}(x)&\equiv
   F_{\mu\rho}^a(x)F_{\nu\rho}^a(x),
\label{eq:(4.6)}\\
   \mathcal{O}_{2\mu\nu}(x)&\equiv
   \delta_{\mu\nu}F_{\rho\sigma}^a(x)F_{\rho\sigma}^a(x),
\label{eq:(4.7)}\\
   \mathcal{O}_{3\mu\nu}(x)&\equiv
   \Bar{\psi}(x)
   \left(\gamma_\mu\overleftrightarrow{D}_\nu
   +\gamma_\nu\overleftrightarrow{D}_\mu\right)
   \psi(x),
\label{eq:(4.8)}\\
   \mathcal{O}_{4\mu\nu}(x)&\equiv
   \delta_{\mu\nu}
   \Bar{\psi}(x)
   \overleftrightarrow{\Slash{D}}
   \psi(x),
\label{eq:(4.9)}\\
   \mathcal{O}_{5\mu\nu}(x)&\equiv
   \delta_{\mu\nu}
   m_0\Bar{\psi}(x)
   \psi(x).
\label{eq:(4.10)}
\end{align}
$\mathcal{O}_{1\mu\nu}(x)$ 与~$\mathcal{O}_{2\mu\nu}(x)$ 的质量量纲为~$4$，
而 $\mathcal{O}_{3\mu\nu}(x)$、$\mathcal{O}_{4\mu\nu}(x)$
和~$\mathcal{O}_{5\mu\nu}(x)$ 的为~$D$。

现在考虑式~\eqref{eq:(4.1)}--\eqref{eq:(4.5)} 中流时间~$t$
非常小的情形。由于流时间为~$t$、位置~$x$ 处的流动场是~$x$ 邻域内半径%
$\sim\sqrt{8t}$ 区域中未流动场的组合，%脚注：这源于流方程基本上是扩散方程这一事实。
故算符~\eqref{eq:(4.1)}--\eqref{eq:(4.5)} 在 $t\to0$
极限下可视为 $x$ 空间中的局域算符。由于式~\eqref{eq:(4.6)}--\eqref{eq:(4.10)}
中的裸算符张成量纲为~$4$（$D\to4$）、CP 偶的二阶对称规范不变局域算符的完备集，
对小~$t$ 我们有如下（渐近）展开：
\begin{equation}
   \Tilde{\mathcal{O}}_{i\mu\nu}(t,x)
   =\left\langle\Tilde{\mathcal{O}}_{i\mu\nu}(t,x)\right\rangle
   +\zeta_{ij}(t)
   \left[\mathcal{O}_{j\mu\nu}(x)
   -\left\langle\mathcal{O}_{j\mu\nu}(x)\right\rangle\right]+O(t),
\label{eq:(4.11)}
\end{equation}
其中被省略的项是量纲 $6$（$D\to4$）或更高的算符的贡献。

在写下小流时间展开~\eqref{eq:(4.11)} 时，我们假定了点~$x$ 处没有其他 $D$
维复合算符。当点~$x$ 处还有其他算符［记作 $\mathcal{P}(x)$］时，展开%
~\eqref{eq:(4.11)} 不必成立，理由如下：式~\eqref{eq:(4.11)}
左边与~$\mathcal{P}(x)$ 的乘积在 $t>0$ 时不具有任何发散，而右边的每一项因是 $D$
维（即未流动的）复合算符，可以与~$\mathcal{P}(x)$ 发生等点奇异性。若式%
~\eqref{eq:(4.11)} 成立，这就是矛盾。这表明下面我们将要导出的能动量张量公式
\emph{只在能动量张量与其他算符没有重叠时成立}。%脚注：这一要点在文
献~\cite{Suzuki:2013gza}中未被充分认识。特别地，关于我们的构造是否重现
Ward--Takahashi 关系~\eqref{eq:(2.9)}，我们无法作出任何论断。尽管如此，由于我们的构造的
归一化是通过与维数正规化中满足式~\eqref{eq:(2.9)} 的能动量张量匹配而确定的，它预期具有正
确的归一。因此，我们的格点能动量张量可用于计算能动量张量与其他算符分离的那些关联函数。
例如与黏滞系数相关的关联函数~\cite{Nakamura:2004sy,Meyer:2007ic,Meyer:2007dy}正是如此。

展开~\eqref{eq:(4.11)} 可反解为
\begin{equation}
   \mathcal{O}_{i\mu\nu}(x)-\left\langle\mathcal{O}_{i\mu\nu}(x)\right\rangle
   =\left(\zeta^{-1}\right)_{ij}(t)
   \left[\Tilde{\mathcal{O}}_{j\mu\nu}(t,x)
   -\left\langle\Tilde{\mathcal{O}}_{j\mu\nu}(t,x)\right\rangle\right]
   +O(t),
\label{eq:(4.12)}
\end{equation}
其中 $\zeta^{-1}$ 表示~$\zeta$ 的逆矩阵。于是，把该关系代入以裸算符表出的能动量张量%
~\eqref{eq:(2.7)} 与~\eqref{eq:(2.10)}：
\begin{align}
   \left\{T_{\mu\nu}\right\}_R(x)
   &=\frac{1}{g_0^2}\left\{
   \mathcal{O}_{1\mu\nu}(x)
   -\left\langle\mathcal{O}_{1\mu\nu}(x)\right\rangle
   -\frac{1}{4}
   \left[\mathcal{O}_{2\mu\nu}(x)
   -\left\langle\mathcal{O}_{2\mu\nu}(x)\right\rangle\right]
   \right\}
\notag\\
   &\qquad{}
   +\frac{1}{4}
   \left[\mathcal{O}_{3\mu\nu}(x)
   -\left\langle\mathcal{O}_{3\mu\nu}(x)\right\rangle\right]
   -\frac{1}{2}
   \left[\mathcal{O}_{4\mu\nu}(x)
   -\left\langle\mathcal{O}_{4\mu\nu}(x)\right\rangle\right]
\notag\\
   &\qquad\qquad{}
   -\left[\mathcal{O}_{5\mu\nu}(x)
   -\left\langle\mathcal{O}_{5\mu\nu}(x)\right\rangle\right],
\label{eq:(4.13)}
\end{align}
便得到表达式（对 $D=4$）
\begin{align}
   \left\{T_{\mu\nu}\right\}_R(x)
   &=c_1(t)\left[
   \Tilde{\mathcal{O}}_{1\mu\nu}(t,x)
   -\frac{1}{4}\Tilde{\mathcal{O}}_{2\mu\nu}(t,x)
   \right]
\notag\\
   &\qquad{}
   +c_2(t)\left[
   \Tilde{\mathcal{O}}_{2\mu\nu}(t,x)
   -\left\langle\Tilde{\mathcal{O}}_{2\mu\nu}(t,x)\right\rangle
   \right]
\notag\\
   &\qquad\qquad{}
   +c_3(t)\left[
   \Tilde{\mathcal{O}}_{3\mu\nu}(t,x)
   -2\Tilde{\mathcal{O}}_{4\mu\nu}(t,x)
   -\left\langle
   \Tilde{\mathcal{O}}_{3\mu\nu}(t,x)
   -2\Tilde{\mathcal{O}}_{4\mu\nu}(t,x)
   \right\rangle
   \right]
\notag\\
   &\qquad\qquad\qquad{}
   +c_4(t)\left[
   \Tilde{\mathcal{O}}_{4\mu\nu}(t,x)
   -\left\langle\Tilde{\mathcal{O}}_{4\mu\nu}(t,x)\right\rangle
   \right]
\notag\\
   &\qquad\qquad\qquad\qquad{}
   +c_5(t)\left[
   \Tilde{\mathcal{O}}_{5\mu\nu}(t,x)
   -\left\langle\Tilde{\mathcal{O}}_{5\mu\nu}(t,x)\right\rangle
   \right]+O(t),
\label{eq:(4.14)}
\end{align}
其中
\begin{align}
   &c_1(t)=\Tilde{c}_1(t),\qquad
   c_2(t)=\Tilde{c}_2(t)+\frac{1}{4}c_1(t),
\notag\\
   &c_3(t)=\Tilde{c}_3(t),\qquad
   c_4(t)=\Tilde{c}_4(t)+2c_3(t),\qquad
   c_5(t)=\Tilde{c}_5(t),
\label{eq:(4.15)}
\end{align}
且
\begin{equation}
   \Tilde{c}_i(t)\equiv\frac{1}{g_0^2}\left\{
   \left(\zeta^{-1}\right)_{1i}(t)
   -\frac{1}{4}\left(\zeta^{-1}\right)_{2i}(t)\right\}
   +\frac{1}{4}\left(\zeta^{-1}\right)_{3i}(t)
   -\frac{1}{2}\left(\zeta^{-1}\right)_{4i}(t)
   -\left(\zeta^{-1}\right)_{5i}(t).
\label{eq:(4.16)}
\end{equation}
在式~\eqref{eq:(4.14)} 中，我们利用了如下事实：有限算
符 $\Tilde{\mathcal{O}}_{1\mu\nu}(t,x)-(1/4)\Tilde{\mathcal{O}}_{2\mu\nu}(t,x)$
在~$D=4$ 无迹，因而没有真空期望值。
式~\eqref{eq:(4.14)} 表明，只要知道系数~$c_i(t)$ 的 $t\to0$ 行为，能动量张量便可作为右边
组合的 $t\to0$ 极限得到。如前所述，既然由（加圈的）流动场构造的复合算符%
~\eqref{eq:(4.1)}--\eqref{eq:(4.5)} 应当与所采用的正规化无关，就可以用格点正规化来计算
式~\eqref{eq:(4.14)} 右边量的关联函数。这提供了用格点正规化计算正确归一且守恒的能动量张
量关联函数的一种可能方法。

于是我们关心式~\eqref{eq:(4.14)} 中系数~$c_i(t)$ 的 $t\to0$ 行为。十分有趣的是，得益于渐
近自由，可以论证系数~$c_i(t)$ 在 $t\to0$ 时可用微扰论求值。为看清这一点，我们对式%
~\eqref{eq:(4.14)} 两边施加
\begin{equation}
   \left(\mu\frac{\partial}{\partial\mu}\right)_0
\label{eq:(4.17)}
\end{equation}
其中 $\mu$ 是重正化标度，下标~$0$ 表示求导时所有裸量保持固定。由于能动量张量不接受相乘重
正化，
$(\mu\partial/\partial\mu)_0(\text{式~\eqref{eq:(4.14)} 的左边})=0$。
而在右边，由于 $\Tilde{O}_{1,2,3,4\mu\nu}(t,x)$
和~$(1/m)\Tilde{O}_{5\mu\nu}(t,x)$ 完全由经流方程的裸量表出，我们有
\begin{equation}
   \left(\mu\frac{\partial}{\partial\mu}\right)_0\Tilde{O}_{1,2,3,4\mu\nu}(t,x)
   =\left(\mu\frac{\partial}{\partial\mu}\right)_0
   \frac{1}{m}\Tilde{O}_{5\mu\nu}(t,x)
   =0.
\label{eq:(4.18)}
\end{equation}
这些观察意味着
\begin{equation}
   \left(\mu\frac{\partial}{\partial\mu}\right)_0c_{1,2,3,4}(t)
   =\left(\mu\frac{\partial}{\partial\mu}\right)_0mc_5(t)
   =0.
\label{eq:(4.19)}
\end{equation}
于是标准重正化群论证说：若把这些量中的重正化参数换成由
\begin{align}
   &q\frac{\mathrm{d}\Bar{g}(q)}{\mathrm{d}q}=\beta(\Bar{g}(q)),\qquad
   \Bar{g}(q=\mu)=g,
\label{eq:(4.20)}\\
   &q\frac{\mathrm{d}\Bar{m}(q)}{\mathrm{d}q}=-\gamma_m(\Bar{g}(q))\Bar{m}(q),\qquad
   \Bar{m}(q=\mu)=m,
\label{eq:(4.21)}
\end{align}
定义的跑动参数（$\mu$ 为原重正化标度），则 $c_{1,2,3,4}(t)$
和~$mc_5(t)$ 与重正化标度无关。这样，既然 $c_{1,2,3,4}(t)$
和~$mc_5(t)$ 与重正化标度无关，两种选择 $q=\mu$ 与~$q=1/\sqrt{8t}$
就应给出相同的结果。由此我们推断
\begin{align}
   c_{1,2,3,4}(t)(g,m;\mu)
   &=c_{1,2,3,4}(t)(\Bar{g}(1/\sqrt{8t}),\Bar{m}(1/\sqrt{8t});1/\sqrt{8t}),
\label{eq:(4.22)}\\
   c_5(t)(g,m;\mu)
   &=\frac{\Bar{m}(1/\sqrt{8t})}{m}
   c_5(t)(\Bar{g}(1/\sqrt{8t}),\Bar{m}(1/\sqrt{8t});1/\sqrt{8t}),
\label{eq:(4.23)}
\end{align}
其中我们显式写出了 $c_i(t)$ 对重正化参数及重正化标度的依赖。最后，得益于渐近自由，跑动规
范耦合 $\Bar{g}(1/\sqrt{8t})\to0$（$t\to0$），因此预期可以用微扰论计算 $t\to0$
的 $c_i(t)$。尽管我们感兴趣的是微扰论无能为力的低能物理，
但式~\eqref{eq:(4.14)} 中 $t\to0$ 的系数~$c_i(t)$ 可以用微扰论求值；这或许可视为某种意义
上的因子化。

\subsection{$c_i(t)$ 至单圈阶}
于是我们在单圈近似下求上式即式~\eqref{eq:(4.14)} 的系数~$c_i(t)$。为此，先把式%
~\eqref{eq:(4.11)} 中的混合系数 $\zeta_{ij}(t)$ 计算到单圈阶；然后 $c_i(t)$
由式~\eqref{eq:(4.15)} 和~\eqref{eq:(4.16)} 得到。$\zeta_{ij}(t)$ 的圈展开应给出
\begin{equation}
   \zeta_{ij}(t)=\delta_{ij}+\zeta_{ij}^{(1)}(t)+\zeta_{ij}^{(2)}(t)+\dotsb,
\label{eq:(4.24)}
\end{equation}
其中对 $j=1$ 和~$j=2$ 上标记圈阶；对 $j=3$、$4$、和~$5$，计入式~\eqref{eq:(3.26)}
中的因子后取
\begin{equation}
   \zeta_{ij}(t)=Z(\epsilon)
   \left[\delta_{ij}+\zeta_{ij}^{(1)}(t)+\zeta_{ij}^{(2)}(t)+\dotsb\right].
\label{eq:(4.25)}
\end{equation}

如同文献~\cite{Suzuki:2013gza}，$\zeta_{ij}^{(1)}(t)$ 可以直接计算。\footnote{本论文先前
版本未能很好解释如下计算方案的合理性：人们或许会问，为何单圈匹配系
数~$\zeta_{ij}^{(1)}(t)$ 可以仅从关联函数~\eqref{eq:(4.26)} 和~\eqref{eq:(4.29)}
读出，而无须计算把流动复合算
符~$\Tilde{\mathcal{O}}_{i\mu\nu}(t,x)$ 换成裸算符~$\mathcal{O}_{i\mu\nu}(x)$
后的相应关联函数。该合理性直接关系到我们处理红外（IR）发散的方式，我们在附
录~\ref{sec:D} 中补充了详细解释。我们非常感谢本刊一位审稿人在此问题上的建议。}例如，为得
到 $j=1$ 和~$2$ 的 $\zeta_{ij}^{(1)}(t)$，考虑关联函数
\begin{equation}
   \left\langle\Tilde{\mathcal{O}}_{i\mu\nu}(t,x)
   A_\beta^b(y)A_\gamma^c(z)\right\rangle.
\label{eq:(4.26)}
\end{equation}
在我们全文采用的费曼规范中，其结构为
\begin{equation}
   g_0^2\int_k\frac{\mathrm{e}^{ik(x-y)}}{k^2}g_0^2\int_\ell\frac{\mathrm{e}^{i\ell(x-z)}}{\ell^2}
   \delta^{bc}\mathcal{M}_{\mu\nu,\beta\gamma}(k,\ell).
\label{eq:(4.27)}
\end{equation}
把 $\mathcal{M}_{\mu\nu,\beta\gamma}(k,\ell)$ 展开到~$O(k,\ell)$ 后，便可利用下列对应关
系读出算符混合：
\begin{equation}
   ik_\mu i\ell_\nu\delta_{\beta\gamma}\to F_{\mu\rho}^a(x)F_{\nu\rho}^a(x),\qquad
   ik\cdot i\ell\delta_{\mu\nu}\delta_{\beta\gamma}
   \to\frac{1}{4}\delta_{\mu\nu}F_{\rho\sigma}^a(x)F_{\rho\sigma}^a(x).
\label{eq:(4.28)}
\end{equation}
这样就得到 $j=1$ 和~$2$ 的 $\zeta_{ij}^{(1)}(t)$。\footnote{动量积分采用附
录~\ref{sec:B} 中的积分公式。}

类似地，为得到 $j=3$、$4$ 和~$5$ 的 $\zeta_{ij}^{(1)}(t)$，考虑
\begin{equation}
   \left\langle\Tilde{\mathcal{O}}_{i\mu\nu}(t,x)
   \psi(y)\Bar{\psi}(z)\right\rangle,
\label{eq:(4.29)}
\end{equation}
其一般结构为
\begin{equation}
   \int_k\frac{\mathrm{e}^{ik(y-x)}}{i\Slash{k}+m_0}
   \mathcal{M}_{\mu\nu}(k,\ell)
   \int_\ell\frac{\mathrm{e}^{i\ell(x-z)}}{i\Slash{\ell}+m_0}.
\label{eq:(4.30)}
\end{equation}
然后把 $\mathcal{M}_{\mu\nu}(k,\ell)$ 展开到~$O(k)$ 和~$O(\ell)$，并用对应关系
\begin{equation}
   \gamma_\mu i(k+\ell)_\nu+\gamma_\nu i(k+\ell)_\mu\to
   \Bar{\psi}(x)\left[
   \gamma_\mu\overleftrightarrow{D}_\nu+\gamma_\nu\overleftrightarrow{D}_\mu
   \right]\psi(x),\qquad
   \delta_{\mu\nu}\to\delta_{\mu\nu}\Bar{\psi}(x)\psi(x),
\label{eq:(4.31)}
\end{equation}
读出算符混合。

\begin{figure}
\begin{minipage}{0.3\textwidth}
\begin{center}
\includegraphics[width=3cm]{images/A03}
\caption{A03}
\label{fig:14}
\end{center}
\end{minipage}
\begin{minipage}{0.3\textwidth}
\begin{center}
\includegraphics[width=3cm]{images/A04}
\caption{A04}
\label{fig:15}
\end{center}
\end{minipage}
\begin{minipage}{0.3\textwidth}
\begin{center}
\includegraphics[width=3cm]{images/A05}
\caption{A05}
\label{fig:16}
\end{center}
\end{minipage}
\end{figure}

\begin{figure}
\begin{minipage}{0.3\textwidth}
\begin{center}
\includegraphics[width=3cm]{images/A06}
\caption{A06}
\label{fig:17}
\end{center}
\end{minipage}
\begin{minipage}{0.3\textwidth}
\begin{center}
\includegraphics[width=3cm]{images/A07}
\caption{A07}
\label{fig:18}
\end{center}
\end{minipage}
\begin{minipage}{0.3\textwidth}
\begin{center}
\includegraphics[width=3cm]{images/A08}
\caption{A08}
\label{fig:19}
\end{center}
\end{minipage}
\end{figure}

\begin{figure}
\begin{minipage}{0.3\textwidth}
\begin{center}
\includegraphics[width=3cm]{images/A09}
\caption{A09}
\label{fig:20}
\end{center}
\end{minipage}
\begin{minipage}{0.3\textwidth}
\begin{center}
\includegraphics[width=3cm]{images/A10}
\caption{A10}
\label{fig:21}
\end{center}
\end{minipage}
\begin{minipage}{0.3\textwidth}
\begin{center}
\includegraphics[width=3cm]{images/A11}
\caption{A11}
\label{fig:22}
\end{center}
\end{minipage}
\end{figure}

\begin{figure}
\begin{minipage}{0.3\textwidth}
\begin{center}
\includegraphics[width=3cm]{images/A12}
\caption{A12}
\label{fig:23}
\end{center}
\end{minipage}
\begin{minipage}{0.3\textwidth}
\begin{center}
\includegraphics[width=3cm]{images/A13}
\caption{A13}
\label{fig:24}
\end{center}
\end{minipage}
\begin{minipage}{0.3\textwidth}
\begin{center}
\includegraphics[width=3cm]{images/A14}
\caption{A14}
\label{fig:25}
\end{center}
\end{minipage}
\end{figure}

\begin{figure}
\begin{minipage}{0.3\textwidth}
\begin{center}
\includegraphics[width=3cm]{images/A15}
\caption{A15}
\label{fig:26}
\end{center}
\end{minipage}
\begin{minipage}{0.3\textwidth}
\begin{center}
\includegraphics[width=3cm]{images/A16}
\caption{A16}
\label{fig:27}
\end{center}
\end{minipage}
\begin{minipage}{0.3\textwidth}
\begin{center}
\includegraphics[width=3cm]{images/A17}
\caption{A17}
\label{fig:28}
\end{center}
\end{minipage}
\end{figure}

\begin{figure}
\begin{minipage}{0.3\textwidth}
\begin{center}
\includegraphics[width=3cm]{images/A18}
\caption{A18}
\label{fig:29}
\end{center}
\end{minipage}
\begin{minipage}{0.3\textwidth}
\begin{center}
\includegraphics[width=3cm]{images/A19}
\caption{A19}
\label{fig:30}
\end{center}
\end{minipage}
\end{figure}
对 $\zeta_{1j}^{(1)}(t)$，图~\ref{fig:14}--\ref{fig:30} 中的图 A03、A04、A05、A06、A07、
A08、A09、A11、A13、A14、A15、A16 以及 A19 有贡献。\footnote{这些图中虚线表示鬼传播子；
图 A10、A12、A17 与 A18 对应常规的波函数重正化，在计算算符混合时应予剔除。}在这些图及下
面的图中，叉号一般代表式~\eqref{eq:(4.1)}--\eqref{eq:(4.5)}
中的某个复合算符［在本情形中即 $\Tilde{\mathcal{O}}_{1\mu\nu}(t,x)$］。除 A19
外，这些图的结果可以借用文献~\cite{Suzuki:2013gza}。为完整起见，按本文约定将这些结果复录
于附录~\ref{sec:C} 的表~\ref{table:0}。结合图~A19 的贡献，我们得
\begin{align}
   \zeta_{11}^{(1)}(t)
   &=\frac{g_0^2}{(4\pi)^2}C_2(G)
   \left[\frac{11}{3}\epsilon(t)^{-1}+\frac{7}{3}\right],
\label{eq:(4.32)}\\
   \zeta_{12}^{(1)}(t)
   &=\frac{g_0^2}{(4\pi)^2}C_2(G)
   \left[-\frac{11}{12}\epsilon(t)^{-1}-\frac{1}{6}\right],
\label{eq:(4.33)}\\
   \zeta_{13}^{(1)}(t)
   &=\frac{g_0^4}{(4\pi)^2}C_2(R)
   \left[-\frac{2}{3}\epsilon(t)^{-1}-\frac{7}{18}\right],
\label{eq:(4.34)}\\
   \zeta_{14}^{(1)}(t)
   &=\frac{g_0^4}{(4\pi)^2}C_2(R)
   \left[\frac{1}{3}\epsilon(t)^{-1}-\frac{7}{18}\right],
\label{eq:(4.35)}\\
   \zeta_{15}^{(1)}(t)
   &=\frac{g_0^4}{(4\pi)^2}C_2(R)
   \left[3\epsilon(t)^{-1}+\frac{5}{2}\right],
\label{eq:(4.36)}
\end{align}
其中
\begin{equation}
   \epsilon(t)^{-1}\equiv\frac{1}{\epsilon}+\ln(8\pi t).
\label{eq:(4.37)}
\end{equation}
考虑 $\Tilde{\mathcal{O}}_{1\mu\nu}$ 的迹部分，由这些进一步得到
\begin{align}
   \zeta_{21}^{(1)}(t)
   &=0,
\label{eq:(4.38)}\\
   \zeta_{22}^{(1)}(t)
   &=\zeta_{11}^{(1)}(t)+\zeta_{12}^{(1)}(t)(4-2\epsilon)
   =\frac{g_0^2}{(4\pi)^2}C_2(G)\frac{7}{2},
\label{eq:(4.39)}\\
   \zeta_{23}^{(1)}(t)
   &=0,
\label{eq:(4.40)}\\
   \zeta_{24}^{(1)}(t)
   &=2\zeta_{13}^{(1)}(t)+\zeta_{14}^{(1)}(t)(4-2\epsilon)
   =\frac{g_0^4}{(4\pi)^2}C_2(R)(-3),
\label{eq:(4.41)}\\
   \zeta_{25}^{(1)}(t)
   &=\zeta_{15}^{(1)}(t)(4-2\epsilon)
   =\frac{g_0^4}{(4\pi)^2}C_2(R)
   \left[12\epsilon(t)^{-1}+4\right].
\label{eq:(4.42)}
\end{align}

\begin{figure}
\begin{minipage}{0.3\textwidth}
\begin{center}
\includegraphics[width=3cm]{images/B03}
\caption{B03}
\label{fig:31}
\end{center}
\end{minipage}
\begin{minipage}{0.3\textwidth}
\begin{center}
\includegraphics[width=3cm]{images/B04}
\caption{B04}
\label{fig:32}
\end{center}
\end{minipage}
\begin{minipage}{0.3\textwidth}
\begin{center}
\includegraphics[width=3cm]{images/B05}
\caption{B05}
\label{fig:33}
\end{center}
\end{minipage}
\end{figure}

\begin{figure}
\begin{minipage}{0.3\textwidth}
\begin{center}
\includegraphics[width=3cm]{images/B06}
\caption{B06}
\label{fig:34}
\end{center}
\end{minipage}
\begin{minipage}{0.3\textwidth}
\begin{center}
\includegraphics[width=3cm]{images/B07}
\caption{B07}
\label{fig:35}
\end{center}
\end{minipage}
\begin{minipage}{0.3\textwidth}
\begin{center}
\includegraphics[width=3cm]{images/B08}
\caption{B08}
\label{fig:36}
\end{center}
\end{minipage}
\end{figure}

\begin{figure}
\begin{minipage}{0.3\textwidth}
\begin{center}
\includegraphics[width=3cm]{images/B09}
\caption{B09}
\label{fig:37}
\end{center}
\end{minipage}
\begin{minipage}{0.3\textwidth}
\begin{center}
\includegraphics[width=3cm]{images/B10}
\caption{B10}
\label{fig:38}
\end{center}
\end{minipage}
\begin{minipage}{0.3\textwidth}
\begin{center}
\includegraphics[width=3cm]{images/B11}
\caption{B11}
\label{fig:39}
\end{center}
\end{minipage}
\end{figure}

\begin{figure}
\begin{minipage}{0.3\textwidth}
\begin{center}
\includegraphics[width=3cm]{images/B12}
\caption{B12}
\label{fig:40}
\end{center}
\end{minipage}
\begin{minipage}{0.3\textwidth}
\begin{center}
\includegraphics[width=3cm]{images/B13}
\caption{B13}
\label{fig:41}
\end{center}
\end{minipage}
\begin{minipage}{0.3\textwidth}
\begin{center}
\includegraphics[width=3cm]{images/B14}
\caption{B14}
\label{fig:42}
\end{center}
\end{minipage}
\end{figure}

\begin{figure}
\begin{minipage}{0.3\textwidth}
\begin{center}
\includegraphics[width=3cm]{images/B15}
\caption{B15}
\label{fig:43}
\end{center}
\end{minipage}
\begin{minipage}{0.3\textwidth}
\begin{center}
\includegraphics[width=3cm]{images/B16}
\caption{B16}
\label{fig:44}
\end{center}
\end{minipage}
\begin{minipage}{0.3\textwidth}
\begin{center}
\includegraphics[width=3cm]{images/B17}
\caption{B17}
\label{fig:45}
\end{center}
\end{minipage}
\end{figure}

\begin{figure}
\begin{minipage}{0.3\textwidth}
\begin{center}
\includegraphics[width=3cm]{images/B18}
\caption{B18}
\label{fig:46}
\end{center}
\end{minipage}
\end{figure}

对于 $j=1$ 和~$2$ 的 $\zeta_{3j}^{(1)}(t)$，图~\ref{fig:31}--\ref{fig:40}
中的图 B03、B04、B08、B09、B10、B11 及~B12 有贡献。每个图的贡献列于表~\ref{table:2}。
对于 $j=3$、$4$、和~$5$ 的 $\zeta_{3j}^{(1)}(t)$，图~\ref{fig:34}--\ref{fig:46}
中的图 B06、B07、B13、B14、B15、B16、B17 与 B18 有贡献，其贡献列于表~\ref{table:3}。
%
\begin{table}
\caption{$\zeta_{3j}^{(1)}$，单位为 $\frac{1}{(4\pi)^2}T(R)N_{\mathrm{f}}$}
\label{table:2}
\begin{center}
\renewcommand{\arraystretch}{2.2}
\setlength{\tabcolsep}{20pt}
\begin{tabular}{crr}
\toprule
 diagram & \multicolumn{1}{c}{$\zeta_{31}^{(1)}(t)$}
 & \multicolumn{1}{c}{$\zeta_{32}^{(1)}(t)$} \\
\midrule
 B03  & $-\dfrac{16}{3}\epsilon(t)^{-1}-\dfrac{64}{9}$ & $\dfrac{4}{3}\epsilon(t)^{-1}+\dfrac{25}{9}$ \\
 B04  & $0$ & $0$ \\
 B08  & $0$ & $-\dfrac{1}{3}$ \\
 B09  & $-\dfrac{5}{12}$ & $-\dfrac{23}{144}$ \\
 B10  & $\dfrac{5}{12}$ & $-\dfrac{1}{16}$ \\
 B11  & $\dfrac{10}{9}$ & $\dfrac{7}{9}$ \\
 B12  & $0$ & $0$ \\
\bottomrule
\end{tabular}
\end{center}
\end{table}
%
\begin{table}
\caption{$\zeta_{3j}^{(1)}$，单位为 $\frac{g_0^2}{(4\pi)^2}C_2(R)$}
\label{table:3}
\begin{center}
\renewcommand{\arraystretch}{2.2}
\setlength{\tabcolsep}{20pt}
\begin{tabular}{crrr}
\toprule
 diagram & \multicolumn{1}{c}{$\zeta_{33}^{(1)}(t)$}
 & \multicolumn{1}{c}{$\zeta_{34}^{(1)}(t)$}
 & \multicolumn{1}{c}{$\zeta_{35}^{(1)}(t)$} \\
\midrule
 B06 & $-\dfrac{1}{3}\epsilon(t)^{-1}+\dfrac{1}{18}$ & $\dfrac{2}{3}\epsilon(t)^{-1}+\dfrac{5}{9}$ & $8\epsilon(t)^{-1}+8$ \\
 B07 & $2\epsilon(t)^{-1}+2$ & $-2\epsilon(t)^{-1}-2$ & $-8\epsilon(t)^{-1}-8$ \\
 B13 & $0$ & $1$ & $0$ \\
 B14 & $0$ & $0$ & $0$ \\
 B15 & $2\epsilon(t)^{-1}+2$ & $0$ & $0$ \\
 B16 & $0$ & $0$ & $0$ \\
 B17 & $-2$ & $0$ & $0$ \\
 B18 & $-4\epsilon(t)^{-1}-2$ & $0$ & $0$ \\
\bottomrule
\end{tabular}
\end{center}
\end{table}
%
作为这些贡献之和，我们得
\begin{align}
   \zeta_{31}^{(1)}(t)
   &=\frac{1}{(4\pi)^2}T(R)N_{\mathrm{f}}
   \left[-\frac{16}{3}\epsilon(t)^{-1}-6\right],
\label{eq:(4.43)}\\
   \zeta_{32}^{(1)}(t)
   &=\frac{1}{(4\pi)^2}T(R)N_{\mathrm{f}}
   \left[\frac{4}{3}\epsilon(t)^{-1}+3\right],
\label{eq:(4.44)}\\
   \zeta_{33}^{(1)}(t)
   &=\frac{g^2}{(4\pi)^2}C_2(R)
   \left[3\frac{1}{\epsilon}-\Phi(t)\right]
   +\frac{g_0^2}{(4\pi)^2}C_2(R)
   \left[-\frac{1}{3}\epsilon(t)^{-1}+\frac{1}{18}\right],
\label{eq:(4.45)}\\
   \zeta_{34}^{(1)}(t)
   &=\frac{g_0^2}{(4\pi)^2}C_2(R)
   \left[-\frac{4}{3}\epsilon(t)^{-1}-\frac{4}{9}\right],
\label{eq:(4.46)}\\
   \zeta_{35}^{(1)}(t)
   &=0,
\label{eq:(4.47)}
\end{align}
其中在 $\zeta_{33}^{(1)}(t)$~\eqref{eq:(4.45)}
中，右边的第一项来自式~\eqref{eq:(3.20)} 和~\eqref{eq:(3.21)}
中从非加圈场到加圈场的换基——回想式~\eqref{eq:(3.25)}；组合~$\Phi(t)$ 由
式~\eqref{eq:(3.27)} 给出。由此进一步得
\begin{align}
   \zeta_{41}^{(1)}(t)
   &=0,
\label{eq:(4.48)}\\
   \zeta_{42}^{(1)}(t)
   &=\frac{1}{2}\zeta_{31}^{(1)}(t)+\frac{1}{2}\zeta_{32}^{(1)}(t)(4-2\epsilon)
   =\frac{1}{(4\pi)^2}T(R)N_{\mathrm{f}}
   \frac{5}{3},
\label{eq:(4.49)}\\
   \zeta_{43}^{(1)}(t)
   &=0,
\label{eq:(4.50)}\\
   \zeta_{44}^{(1)}(t)
   &=\zeta_{33}^{(1)}(t)+\frac{1}{2}\zeta_{34}^{(1)}(t)(4-2\epsilon)
\notag\\
   &=\frac{g^2}{(4\pi)^2}C_2(R)
   \left[3\frac{1}{\epsilon}-\Phi(t)\right]
   +\frac{g_0^2}{(4\pi)^2}C_2(R)
   \left[-3\epsilon(t)^{-1}+\frac{1}{2}\right],
\label{eq:(4.51)}\\
   \zeta_{45}^{(1)}(t)
   &=0.
\label{eq:(4.52)}
\end{align}

最后，
\begin{equation}
   \zeta_{51}^{(1)}(t)
   =\zeta_{52}^{(1)}(t)=\zeta_{53}^{(1)}(t)=\zeta_{54}^{(1)}(t)=0,
\label{eq:(4.53)}
\end{equation}
而 $\zeta_{55}^{(1)}(t)$ 由表~\ref{table:4} 中各单圈图的贡献与到加圈场的换基因子之和给
出：
%
\begin{table}
\caption{$\zeta_{55}^{(1)}(t)$，单位为~$\frac{g_0^2}{(4\pi)^2}C_2(R)$。}
\label{table:4}
\begin{center}
\renewcommand{\arraystretch}{2.2}
\setlength{\tabcolsep}{20pt}
\begin{tabular}{cr}
\toprule
 diagram & \multicolumn{1}{c}{$\zeta_{55}^{(1)}(t)$} \\
\midrule
 B06 & $-4\epsilon(t)^{-1}-2$ \\
 B13 & $0$ \\
 B14 & $0$ \\
 B15 & $2\epsilon(t)^{-1}+2$ \\
 B16 & $0$ \\
 B18 & $-4\epsilon(t)^{-1}-2$ \\
\bottomrule
\end{tabular}
\end{center}
\end{table}
%
\begin{align}
   \zeta_{55}^{(1)}(t)
   =\frac{g^2}{(4\pi)^2}C_2(R)
   \left[6\frac{1}{\epsilon}-\Phi(t)\right]
   +\frac{g_0^2}{(4\pi)^2}C_2(R)
   \left[-6\epsilon(t)^{-1}-2\right].
\label{eq:(4.54)}
\end{align}

至此，我们已得到式~\eqref{eq:(4.24)} 和~\eqref{eq:(4.25)}
中的全部 $\zeta_{ij}^{(1)}(t)$。由于树图级近似下的矩阵~$\zeta_{ij}(t)$ 是单位矩阵，单圈
近似下矩阵~$\zeta_{ij}(t)$ 的求逆是直截了当的；
于是式~\eqref{eq:(4.15)} 和~\eqref{eq:(4.16)} 给出
\begin{align}
   c_1(t)&=\frac{1}{g_0^2}
   \left[1-\zeta_{11}^{(1)}(t)\right]
   -\frac{1}{4}\zeta_{31}^{(1)}(t),
\label{eq:(4.55)}\\
   c_2(t)&=\frac{1}{g_0^2}
   \left(-\frac{1}{2}\epsilon\right)\zeta_{12}^{(1)}(t)
   +\left(\frac{3}{4}-\frac{1}{2}\epsilon\right)\zeta_{32}^{(1)}(t)
   +\frac{3}{16}\zeta_{31}^{(1)}(t),
\label{eq:(4.56)}\\
   c_3(t)&=\left\{\frac{1}{g_0^2}
   \left[-\zeta_{13}^{(1)}(t)\right]
   +\frac{1}{4}-\frac{1}{4}\zeta_{33}^{(1)}(t)\right\}Z(\epsilon)^{-1},
\label{eq:(4.57)}\\
   c_4(t)&=\left\{\frac{1}{g_0^2}
   \left[-\frac{1}{2}\epsilon\zeta_{14}^{(1)}(t)
   -\frac{3}{2}\zeta_{13}^{(1)}(t)\right]
   +\left(\frac{3}{4}-\frac{1}{2}\epsilon\right)\zeta_{34}^{(1)}(t)\right\}
   Z(\epsilon)^{-1},
\label{eq:(4.58)}\\
   c_5(t)&=\left\{\frac{1}{g_0^2}
   \left[-\frac{1}{2}\epsilon\zeta_{15}^{(1)}(t)\right]
   -1+\zeta_{55}^{(1)}(t)\right\}Z(\epsilon)^{-1}.
\label{eq:(4.59)}
\end{align}
然后，使用 $\text{MS}$ 方案的重正化规范耦合，即式~\eqref{eq:(A1)}，在单圈阶我们有（%
$\epsilon\to0$）
\begin{align}
   c_1(t)&
   =\frac{1}{g^2}-b_0\ln(8\pi\mu^2t)-\frac{1}{(4\pi)^2}
   \left[\frac{7}{3}C_2(G)
   -\frac{3}{2}T(R)N_{\mathrm{f}}\right],
\label{eq:(4.60)}\\
   c_2(t)&
   =\frac{1}{8}
   \frac{1}{(4\pi)^2}
   \left[\frac{11}{3}C_2(G)
   +\frac{11}{3}T(R)N_{\mathrm{f}}\right],
\label{eq:(4.61)}\\
   c_3(t)&
   =\frac{1}{4}\left\{1+\frac{g^2}{(4\pi)^2}C_2(R)
   \left[\frac{3}{2}+\ln(432)\right]\right\},
\label{eq:(4.62)}\\
   c_4(t)&=\frac{1}{8}d_0g^2,
\label{eq:(4.63)}\\
   c_5(t)&=-\left\{1+\frac{g^2}{(4\pi)^2}C_2(R)
   \left[3\ln(8\pi\mu^2t)+\frac{7}{2}+\ln(432)\right]\right\}.
\label{eq:(4.64)}
\end{align}
由于式~\eqref{eq:(4.14)} 中的 $c_i(t)$ 连接有限的能动量张量与由（加圈的）流动场构造的 UV
有限局域乘积 $\Tilde{\mathcal{O}}_{i\mu\nu}(t,x)$，它们应当是 UV 有限的。我们对~$c_i(t)$
的显式单圈计算确认了这一有限性，令人相当放心。

若偏好 $\overline{\text{MS}}$ 方案而非上面表达式所假定的 $\text{MS}$
方案，只须作替换
\begin{equation}
    \mu^2\to\frac{\mathrm{e}^{\gamma_{\mathrm{E}}}}{4\pi}\mu^2,
\label{eq:(4.65)}
\end{equation}
其中 $\gamma_{\mathrm{E}}$ 是欧拉常数。

\subsection{一致性检验：迹反常}
有趣的是可以看到，带
式~\eqref{eq:(4.60)}--\eqref{eq:(4.64)} 中 $c_i(t)$ 的式~\eqref{eq:(4.14)}
实际上在单圈近似下重现迹反常~\eqref{eq:(2.11)}。乍看之下，
式~\eqref{eq:(4.61)} 中的 $c_2(t)$ 似乎与正确的迹反常不相容，因为对~$D=4$
有 $\delta_{\mu\nu}[\Tilde{\mathcal{O}}_{1\mu\nu}
-(1/4)\Tilde{\mathcal{O}}_{2\mu\nu}]=0$
且 $\delta_{\mu\nu}\Tilde{\mathcal{O}}_{2\mu\nu}
=4\left\{F_{\rho\sigma}^aF_{\rho\sigma}^a\right\}_R(x)+O(g^2)$［最后的等式来自式%
~\eqref{eq:(A17)}］；另一方面，来自式~\eqref{eq:(4.61)} 的 $4c_2(t)$
并不等同于~$b_0/2$——$b_0$ 是 $\beta$ 函数~\eqref{eq:(2.16)}
的首项系数，也是迹反常的正确单圈系数。

然而这是一个过早的判断。事实上，式~\eqref{eq:(4.49)} 中的 $\zeta_{42}^{(1)}(t)$
表明存在形如
\begin{equation}
   \mathring{\Bar{\chi}}(t,x)
   \overleftrightarrow{\Slash{D}}
   \mathring{\chi}(t,x)
   =\left\{\Bar{\psi}\overleftrightarrow{\Slash{D}}\psi\right\}_R(x)
   +\frac{1}{(4\pi)^2}T(R)N_{\mathrm{f}}\frac{5}{3}
   \left\{F_{\rho\sigma}^aF_{\rho\sigma}^a\right\}_R(x)
   +\dotsb.
\label{eq:(4.66)}
\end{equation}
的算符混合。于是最后一项恰好填补了 $4c_2(t)$
与~$b_0/2$ 之间的差。

这样一来，在单圈阶我们有
\begin{align}
   \delta_{\mu\nu}\left\{T_{\mu\nu}\right\}_R(x)
   &=\frac{1}{2}\frac{1}{(4\pi)^2}
   \left[\frac{11}{3}C_2(G)-\frac{4}{3}T(R)N_{\mathrm{f}}\right]
   \left\{F_{\rho\sigma}^aF_{\rho\sigma}^a\right\}_R(x)
\notag\\
   &\qquad{}
   -\frac{3}{2}
   \left\{\Bar{\psi}\overleftrightarrow{\Slash{D}}\psi\right\}_R(x)
   -\left[4+\frac{g^2}{(4\pi)^2}6C_2(R)\right]
   m\left\{\Bar{\psi}\psi\right\}_R(x).
\label{eq:(4.67)}
\end{align}
若使用重正化场的运动方程
\begin{equation}
   \left\{\Bar{\psi}\overleftrightarrow{\Slash{D}}\psi\right\}_R(x)
   =-2m\left\{\Bar{\psi}\psi\right\}_R(x),
\label{eq:(4.68)}
\end{equation}
——正如我们所假设的点~$x$ 处无其他算符时其使用是正当的——则它在单圈层次重现迹反常%
~\eqref{eq:(2.11)}。我们看到，式~\eqref{eq:(4.60)}--\eqref{eq:(4.64)}
中的单圈结果与迹反常一致。

在文献~\cite{Suzuki:2013gza}中，针对纯 Yang--Mills 理论，通过要求表达
式~\eqref{eq:(4.14)} 在双圈层次重现迹反常~\eqref{eq:(2.11)}，确定了 $c_2(t)$
的次领头（双圈）项为
\begin{equation}
   c_2(t)=\frac{1}{8}b_0
   -\frac{1}{8}b_0
   \left[\frac{1}{(4\pi)^2}G_2(G)\frac{7}{2}-\frac{b_1}{b_0}\right]g^2,
\label{eq:(4.69)}
\end{equation}
其中 $b_0=[1/(4\pi)^2](11/3)C_2(G)$ 及~$b_1=[1/(4\pi)^4](34/3)C_2(G)^2$。但对含费米子的
现系统而言，仅凭这一要求似乎无法同时固定 $c_2(t)$ 与 $c_4(t)$
中的次领头项；因此在本文处理含费米子的系统时，我们满足于单圈公式，
即式~\eqref{eq:(4.60)}--\eqref{eq:(4.64)}。

\subsection{主公式}
由式~\eqref{eq:(4.14)}，能动量张量由 $t\to0$ 极限给出：
\begin{align}
   \left\{T_{\mu\nu}\right\}_R(x)
   &=\lim_{t\to0}\biggl\{c_1(t)\left[
   \Tilde{\mathcal{O}}_{1\mu\nu}(t,x)
   -\frac{1}{4}\Tilde{\mathcal{O}}_{2\mu\nu}(t,x)
   \right]
\notag\\
   &\qquad\qquad{}
   +c_2(t)\left[
   \Tilde{\mathcal{O}}_{2\mu\nu}(t,x)
   -\left\langle\Tilde{\mathcal{O}}_{2\mu\nu}(t,x)\right\rangle
   \right]
\notag\\
   &\qquad\qquad\qquad{}
   +c_3(t)\left[
   \Tilde{\mathcal{O}}_{3\mu\nu}(t,x)
   -2\Tilde{\mathcal{O}}_{4\mu\nu}(t,x)
   -\left\langle
   \Tilde{\mathcal{O}}_{3\mu\nu}(t,x)
   -2\Tilde{\mathcal{O}}_{4\mu\nu}(t,x)
   \right\rangle
   \right]
\notag\\
   &\qquad\qquad\qquad\qquad{}
   +c_4(t)\left[
   \Tilde{\mathcal{O}}_{4\mu\nu}(t,x)
   -\left\langle\Tilde{\mathcal{O}}_{4\mu\nu}(t,x)\right\rangle
   \right]
\notag\\
   &\qquad\qquad\qquad\qquad\qquad{}
   +c_5(t)\left[
   \Tilde{\mathcal{O}}_{5\mu\nu}(t,x)
   -\left\langle\Tilde{\mathcal{O}}_{5\mu\nu}(t,x)\right\rangle
   \right]\biggr\},
\label{eq:(4.70)}
\end{align}
其中右边的算符由式~\eqref{eq:(4.1)}--\eqref{eq:(4.5)} 给出。还可以进一步利用恒等式
\begin{equation}
   2\left\langle\Tilde{\mathcal{O}}_{3\mu\nu}(t,x)\right\rangle
   =\left\langle\Tilde{\mathcal{O}}_{4\mu\nu}(t,x)\right\rangle
   =\frac{-2\dim(R)N_{\mathrm{f}}}{(4\pi)^2t^2}\delta_{\mu\nu}
\label{eq:(4.71)}
\end{equation}
使表达式稍微简化。把重正化群论证的结论、即式~\eqref{eq:(4.22)}--\eqref{eq:(4.23)}，
应用于式~\eqref{eq:(4.60)}--\eqref{eq:(4.64)}，我们得
\begin{align}
   c_1(t)&
   =\frac{1}{\Bar{g}(1/\sqrt{8t})^2}
   -b_0\ln\pi-\frac{1}{(4\pi)^2}
   \left[\frac{7}{3}C_2(G)
   -\frac{3}{2}T(R)N_{\mathrm{f}}\right],
\label{eq:(4.72)}\\
   c_2(t)&
   =\frac{1}{8}
   \frac{1}{(4\pi)^2}
   \left[\frac{11}{3}C_2(G)
   +\frac{11}{3}T(R)N_{\mathrm{f}}\right],
\label{eq:(4.73)}\\
   c_3(t)&
   =\frac{1}{4}\left\{1+\frac{\Bar{g}(1/\sqrt{8t})^2}{(4\pi)^2}C_2(R)
   \left[\frac{3}{2}+\ln(432)\right]\right\},
\label{eq:(4.74)}\\
   c_4(t)&=\frac{1}{8}d_0\Bar{g}(1/\sqrt{8t})^2,
\label{eq:(4.75)}\\
   c_5(t)&=-\frac{\Bar{m}(1/\sqrt{8t})}{m}
   \left\{1+\frac{\Bar{g}(1/\sqrt{8t})^2}{(4\pi)^2}C_2(R)
   \left[3\ln\pi+\frac{7}{2}+\ln(432)\right]\right\},
\label{eq:(4.76)}
\end{align}
其中 $\Bar{g}(q)$ 是 $\text{MS}$ 方案的跑动规范耦合。要从 $\text{MS}$
方案转到 $\overline{\text{MS}}$ 方案，只须在上面的表达式中作与
式~\eqref{eq:(4.65)} 相对应的替换
\begin{equation}
   \ln\pi\to\gamma_{\mathrm{E}}-2\ln2,
\label{eq:(4.77)}
\end{equation}
即可。带式~\eqref{eq:(4.72)}--\eqref{eq:(4.76)} 的式~\eqref{eq:(4.70)}
就是我们的主要结果。注意 $c_5(t)$~\eqref{eq:(4.76)}
中的重正化质量参数~$m$ 与 $\Tilde{\mathcal{O}}_{5\mu\nu}(t,x)$~\eqref{eq:(4.5)}
中的那个在式~\eqref{eq:(4.70)} 中是冗余的，因为二者在乘积中相互抵消。
对 $t\to0$ 时 $c_5(t)$ 中的跑动质量参
数~$\Bar{m}(1/\sqrt{8t})$，可以使用关系
\begin{align}
   \Bar{m}(q)&=[2b_0\Bar{g}(q)^2]^{d_0/2b_0}
   \exp\left\{\int_0^{\Bar{g}(q)}\mathrm{d}g\,
   \left[-\frac{\gamma_m(g)}{\beta(g)}-\frac{d_0}{b_0g}\right]\right\}m^\infty
\notag\\
   &=\left(\frac{2}{\ell}\right)^{d_0/2b_0}
   \left[1-\frac{d_0b_1}{2b_0^3\ell}(1+\ln\ell)+\frac{d_1}{2b_0^2\ell}
   +O(\ell^{-2})\right]m^\infty,\qquad\ell\equiv\ln(q^2/\Lambda^2),
\label{eq:(4.78)}
\end{align}
其中 $m^\infty$ 表示重正化群不变质量。对有质量的费米子，$m^\infty$~例如可用文
献~\cite{Capitani:1998mq}中确立的方法确定。对无质量费米子，$m^\infty=0$，
我们可以直接弃去式~\eqref{eq:(4.70)} 的最后一行。
